
\subsection{几何光学}
\textbf{反射定律}：反射角等于入射角，入射光线、反射光线与法线共面。\\
\textbf{折射定律}：满足 $n_1\sin\theta_1=n_2\sin\theta_2$，光在不同介质中传播方向发生改变。\\
\textbf{全反射}：光从光密介质射向光疏介质，且入射角大于临界角时发生全反射。

\begin{formulanotebox}
% TODO: 本节公式表格与公式注意事项待补充
\end{formulanotebox}

\begin{keypointbox}
几何光学作图题核心是“法线+角度定义+光路可逆”，全反射要同时满足“从光密到光疏”且“入射角大于临界角”。
\end{keypointbox}

\begin{casetypebox}[title={\strut\textcolor{white}{\bfseries 常见题型：反射定律}}]
% TODO: 反射作图题型解析待补充
\end{casetypebox}

\begin{casetypebox}[title={\strut\textcolor{white}{\bfseries 常见题型：折射定律}}]
% TODO: 折射作图与折射率题型解析待补充
\end{casetypebox}

\begin{casetypebox}[title={\strut\textcolor{white}{\bfseries 常见题型：全反射}}]
% TODO: 全反射临界条件题型解析待补充
\end{casetypebox}